\subsection{09.12.2025}

\subsubsection{Ziel}

Das Ziel des Tages war es, weitere Lektionen der JOY-PI-Anleitung
anzusehen und durchzuführen.

\subsubsection{Durchgeführte Arbeiten}

Es wurde der Quellcode der Lektion 7 zur Erfassung von Temperatur und
Luftfeuchtigkeit analysiert und ausgeführt, da dieser Sensor für das
Projekt essenziell ist.

Zusätzlich wurde der Quellcode der Lektion 10 zur Steuerung des
LCD-Displays betrachtet, da das Display optional zur Darstellung der
Messdaten verwendet werden kann.

\subsubsection{Problem 1: Ausführung der Lektion 7}

Das Programm \texttt{dht\_11.py} wurde innerhalb der aktivierten
virtuellen Umgebung mit dem Befehl
\texttt{sudo python3 dht\_11.py} ausgeführt.

Dabei trat die Fehlermeldung
\texttt{ModuleNotFoundError: No module named 'dht11'} auf.

Beim erneuten Versuch, das Modul zu installieren, wurde folgende
Meldung angezeigt:

\texttt{Requirement already satisfied: dht11 in ./my-venv/lib/python3.13/site-packages}

\texttt{Requirement already satisfied: RPi.GPIO in ./my-venv/lib/python3.13/site-packages}

\textbf{Lösung}

Nach mehreren erfolglosen Versuchen, das Modul erneut zu installieren,
wurde das Skript ohne \texttt{sudo} mit dem Befehl
\texttt{python3 dht\_11.py} ausgeführt.

Dadurch konnte das Programm erfolgreich gestartet werden.
Das Verwenden von \texttt{sudo} war in diesem Fall nicht notwendig,
da die virtuelle Umgebung bereits die benötigten Berechtigungen und
Bibliotheken bereitstellte.

Vermutlich trat das Problem in vorherigen Lektionen nicht auf,
da die benötigten Module dort zusätzlich systemweit installiert waren.

\subsubsection{Problem 2: Ausführung der Lektion 10}

Beim Ausführen des Skripts \texttt{lcd.py} trat die Fehlermeldung auf,
dass das Modul \texttt{board} nicht gefunden werden konnte.

\textbf{Lösung}

An diesem Tag konnte noch keine endgültige Lösung gefunden werden.
Es wurde versucht, das fehlende Modul zu installieren und die Ursache
des Problems weiter zu analysieren.

\subsubsection{Erkenntnisse}

Erkenntnisse aus Lektion 7:

\begin{itemize}
    \item Für die Nutzung des Sensors muss das Modul
    \texttt{dht11} importiert werden.

    \item Das \texttt{dht11}-Objekt wird mit der entsprechenden
    GPIO-Pin-Nummer initialisiert.

    \item Die Methode \texttt{read()} gibt ein Objekt zurück,
    das die Messwerte für Temperatur und Luftfeuchtigkeit enthält.

    \item Es kann zu ungültigen Messwerten kommen, die mit
    \texttt{is\_valid()} überprüft werden können.
\end{itemize}

Erkenntnisse aus Lektion 10:

\begin{itemize}
    \item Die Zeile
    \texttt{\# -*- coding: utf-8 -*-}
    ermöglicht die korrekte Darstellung von Umlauten.

    \item Für die Steuerung des LCD-Displays müssen die Module
    \texttt{time}, \texttt{board}, \texttt{busio} und
    \texttt{adafruit\_character\_lcd.character\_lcd\_i2c}
    importiert werden.

    \item Das LCD-Display muss mit den entsprechenden Pins sowie
    der Anzahl an Zeilen und Spalten initialisiert werden.

    \item Der I\textsuperscript{2}C-Bus wird mit
    \texttt{busio.I2C()} initialisiert.

    \item Mit \texttt{lcd.message}, \texttt{lcd.clear()} und
    \texttt{lcd.backlight} kann das LCD-Display gesteuert werden.
\end{itemize}

\subsubsection{Nächste Schritte}

Für die nächste Woche ist geplant, die Probleme mit Lektion 10 zu lösen,
damit das LCD-Display zur Darstellung der Messdaten verwendet werden kann.